% GENERATED FILE. Edit canonical lessons, then run npm run generate:books.
% canonical-source-hash: fnv1a64:0b3fe0c400ff4d14
% canonical-lessons: RU-C72-here, RU-C72-there, RU-C72-far, RU-C72-near, RU-C72-above

\chapter{Here, There, Far, Near, Up Above}
\label{ch:ru-here-there-far-near-up-above}
\hlchaptermodality{72}

\begin{chapteropening}
\textbf{By the end of this chapter:} \emph{I can say here, there, far, near and up above.}

Complete the last lesson of chapter 72: I can say here, there, far, near and up above.
\end{chapteropening}

\section[\texorpdfstring{\ru{здесь} (zdes)}{здесь (zdes)}]{\ru{здесь} (zdes) — here}

\label{lesson:RU-C72-here}



\begin{quote}
Before the new one: say the Russian for summer, then the Russian for autumn.
\end{quote}

\subsection*{You'll want to know: \ru{здесь}}
\textbf{\ru{здесь}} (\emph{zdes}) — \textquotedblleft{}here\textquotedblright{}.

\begin{cousinweb}
\textbf{\ru{здесь}} $\leftarrow$ Old Russian \emph{\ru{сьде}}, with a later \emph{-\ru{сь}}: Proto-Slavic \emph{*sĭde}, \textquotedblleft{}here\textquotedblright{}, from the old pronoun \emph{*sĭ}, \textquotedblleft{}this\textquotedblright{}, the one in \emph{\ru{сей}} and \emph{\ru{сегодня}}, \textquotedblleft{}today\textquotedblright{}, built as \emph{\ru{где}}, \textquotedblleft{}where\textquotedblright{}, is built from the question word (Vasmer).
\end{cousinweb}

\begin{grammarlens}[title={What you've built}]
One more word for this chapter.
\end{grammarlens}

\subsection*{Your turn}
\begin{itemize}
\raggedright
  \item \emph{Say it:} \emph{zdes}
  \item \emph{Say it:} \emph{zdes}, once more
  \item \emph{Say it:} say \emph{zdes}
  \item \emph{Recall:} say the Russian for summer, then the Russian for autumn, then say \emph{zdes} again
\end{itemize}

\begin{tcolorbox}[breakable,colback=teal!4,colframe=teal!35!black,title={Before you move on}]
What does \ru{здесь} mean? (\textquotedblleft{}Here\textquotedblright{}.) Say it once more.
\end{tcolorbox}
\section[\texorpdfstring{\ru{там} (tam)}{там (tam)}]{\ru{там} (tam) — there}

\label{lesson:RU-C72-there}



\begin{quote}
Before the new one: say the Russian for autumn, then the Russian for here.
\end{quote}

\subsection*{You'll want to know: \ru{там}}
\textbf{\ru{там}} (\emph{tam}) — \textquotedblleft{}there\textquotedblright{}.

\begin{cousinweb}
\textbf{\ru{там}} $\leftarrow$ Old Russian \emph{\ru{тамо}}, Proto-Slavic \emph{*tamo} — inherited, built on the pronoun stem \emph{t-} of \emph{\ru{тот}}, \textquotedblleft{}that\textquotedblright{} (Vasmer).
\end{cousinweb}

\begin{grammarlens}[title={What you've built}]
One more word for this chapter.
\end{grammarlens}

\subsection*{Your turn}
\begin{itemize}
\raggedright
  \item \emph{Say it:} \emph{tam}
  \item \emph{Say it:} \emph{tam}, once more
  \item \emph{Say it:} say \emph{tam}
  \item \emph{Recall:} say the Russian for autumn, then the Russian for here, then say \emph{tam} again
\end{itemize}

\begin{tcolorbox}[breakable,colback=teal!4,colframe=teal!35!black,title={Before you move on}]
What does \ru{там} mean? (\textquotedblleft{}There\textquotedblright{}.) Say it once more.
\end{tcolorbox}
\section[\texorpdfstring{\ru{далеко} (dalekó)}{далеко (dalekó)}]{\ru{далеко} (dalekó) — far}

\label{lesson:RU-C72-far}



\begin{quote}
Before the new one: say the Russian for here, then the Russian for there.
\end{quote}

\subsection*{You'll want to know: \ru{далеко}}
\textbf{\ru{далеко}} (\emph{dalekó}) — \textquotedblleft{}far\textquotedblright{}.

\begin{cousinweb}
\textbf{\ru{далеко}} is \emph{\ru{далёкий}}, \textquotedblleft{}distant\textquotedblright{}, made into an adverb with \emph{-\ru{о}}; it is built on \emph{\ru{даль}}, \textquotedblleft{}the distance\textquotedblright{}, which Vasmer reports is usually linked with \emph{\ru{долгий}}, \textquotedblleft{}long\textquotedblright{}, and Greek \emph{dolichós}, \textquotedblleft{}long\textquotedblright{}.
\end{cousinweb}

\begin{grammarlens}[title={What you've built}]
One more word for this chapter.
\end{grammarlens}

\subsection*{Your turn}
\begin{itemize}
\raggedright
  \item \emph{Say it:} \emph{dalekó}
  \item \emph{Say it:} \emph{dalekó}, once more
  \item \emph{Say it:} say \emph{dalekó}
  \item \emph{Recall:} say the Russian for here, then the Russian for there, then say \emph{dalekó} again
\end{itemize}

\begin{tcolorbox}[breakable,colback=teal!4,colframe=teal!35!black,title={Before you move on}]
What does \ru{далеко} mean? (\textquotedblleft{}Far\textquotedblright{}.) Say it once more.
\end{tcolorbox}
\section[\texorpdfstring{\ru{близко} (blízko)}{близко (blízko)}]{\ru{близко} (blízko) — near, close by}

\label{lesson:RU-C72-near}



\begin{quote}
Before the new one: say the Russian for there, then the Russian for far.
\end{quote}

\subsection*{You'll want to know: \ru{близко}}
\textbf{\ru{близко}} (\emph{blízko}) — \textquotedblleft{}near, close by\textquotedblright{}.

\begin{cousinweb}
\textbf{\ru{близко}} is \emph{\ru{близкий}}, \textquotedblleft{}close\textquotedblright{}, made into an adverb, from \emph{\ru{близ}}, \textquotedblleft{}near\textquotedblright{}, Proto-Slavic \emph{*blizŭ}. Vasmer relates it to Latvian \emph{blaîzît}, \textquotedblleft{}to press\textquotedblright{}, and Latin \emph{flīgere}, \textquotedblleft{}to strike\textquotedblright{}, but also cites a Lithuanian alternative, so the deeper origin is open.
\end{cousinweb}

\begin{grammarlens}[title={What you've built}]
One more word for this chapter.
\end{grammarlens}

\subsection*{Your turn}
\begin{itemize}
\raggedright
  \item \emph{Say it:} \emph{blízko}
  \item \emph{Say it:} \emph{blízko}, once more
  \item \emph{Say it:} say \emph{blízko}
  \item \emph{Recall:} say the Russian for there, then the Russian for far, then say \emph{blízko} again
\end{itemize}

\begin{tcolorbox}[breakable,colback=teal!4,colframe=teal!35!black,title={Before you move on}]
What does \ru{близко} mean? (\textquotedblleft{}Near, close by\textquotedblright{}.) Say it once more.
\end{tcolorbox}
\section[\texorpdfstring{\ru{наверху} (naverkhú)}{наверху (naverkhú)}]{\ru{наверху} (naverkhú) — up above, upstairs}

\label{lesson:RU-C72-above}



\begin{quote}
Before the new one: say the Russian for far, then the Russian for near.
\end{quote}

\subsection*{You'll want to know: \ru{наверху}}
\textbf{\ru{наверху}} (\emph{naverkhú}) — \textquotedblleft{}up above, upstairs\textquotedblright{}.

\begin{cousinweb}
\textbf{\ru{наверху}} is \emph{\ru{на}}, \textquotedblleft{}on\textquotedblright{}, with \emph{\ru{верху}}, a form of \emph{\ru{верх}}, \textquotedblleft{}the top\textquotedblright{}: \textquotedblleft{}on top\textquotedblright{}. \emph{\ru{верх}} $\leftarrow$ Proto-Slavic \emph{*vĭrxŭ}, a cousin of Lithuanian \emph{viršùs}, \textquotedblleft{}top\textquotedblright{} (Vasmer).
\end{cousinweb}

\begin{grammarlens}[title={What you've built}]
One more word for this chapter.
\end{grammarlens}

\subsection*{Your turn}
\begin{itemize}
\raggedright
  \item \emph{Say it:} \emph{naverkhú}
  \item \emph{Say it:} \emph{naverkhú}, once more
  \item \emph{Say it:} say \emph{naverkhú}
  \item \emph{Recall:} say the Russian for far, then the Russian for near, then say \emph{naverkhú} again
\end{itemize}

\begin{tcolorbox}[breakable,colback=teal!4,colframe=teal!35!black,title={Before you move on}]
What does \ru{наверху} mean? (\textquotedblleft{}Up above, upstairs\textquotedblright{}.) Say it once more.
\end{tcolorbox}
